\section{Introduction}


The history of e-commerce platforms follows a pattern of accretion. Systems designed for catalog management acquire shopping carts, then payment processing, then analytics, then---belatedly---AI capabilities. Each addition creates integration complexity while the core architecture, optimized for earlier requirements, constrains what AI can achieve.

Hanzo inverts this pattern. We designed an AI-first commerce platform where machine learning models are not peripheral services but central infrastructure. Every customer interaction---search, browse, cart, checkout---flows through AI systems that personalize, predict, and optimize in real time.

This architectural decision yields compounding benefits:
\begin{enumerate}
\item \textbf{Unified data model}: All events feed into a shared feature store, enabling cross-functional learning.
\item \textbf{Consistent inference}: A single model serving infrastructure ensures low-latency predictions across all surfaces.
\item \textbf{Rapid iteration}: Standardized ML pipelines accelerate model development and deployment.
\item \textbf{Emergent intelligence}: Models learn from each other's predictions, creating feedback loops that improve system-wide performance.
\end{enumerate}

Our contributions are:
\begin{enumerate}
\item An AI-first architecture for commerce platforms with unified feature stores and model serving.
\item Production implementations of personalization, search ranking, fraud detection, and dynamic pricing.
\item Empirical evaluation demonstrating significant improvements across key commerce metrics.
\item Open discussion of challenges and lessons learned in deploying AI at commerce scale.
\end{enumerate}

